\documentclass[a4paper,11pt]{report}
\usepackage{amsmath,amsfonts,amssymb,amsthm,hyperref, marvosym, accents}
%\usepackage{fullpage}
\usepackage{bm}
\usepackage{cleveref}
\usepackage{mdframed}
\usepackage{csquotes}
\usepackage[super]{natbib}
\usepackage{minitoc}
\usepackage{wasysym}
\usepackage{setspace}
\usepackage{enumerate}
\usepackage{enumitem}
\usepackage{tikz-cd}
\usepackage{tikz}

% Optional but nice:
\usepackage[margin=1in]{geometry}

\title{Black boxes, bitter lessons, and compositional semantics}
\author{G.T.} 
\date{\today}

\begin{document}
\maketitle

\section{Motivation}
Sutton's \emph{Bitter Lesson} is the observation that most significant advances in artifical intelligence have come from leveraging increases in processing power rather than the development of novel algorithms to implement the domain-specific knowledge of humans. Wang-Ma\'{s}cianica (2023) claims that the Bitter Lesson applies also to the field of formal semantics  

\section{Hopf algebras of proof trees}
Text...

\subsection{Why decomposition matters}
Text...

\section{Representations and distributional semantics}
Text...

\bibliographystyle{plainnat}
\bibliography{refs}

\end{document}